\documentclass{article}
\usepackage{graphicx}
\usepackage[T1]{fontenc}
\usepackage[utf8]{inputenc}
\usepackage[polish]{babel}
\usepackage[a4paper, margin=1in]{geometry}
\usepackage{amsmath}
\usepackage{hyperref}
\def\figureautorefname{Rys.}
\renewcommand{\equationautorefname}{Równ.}
\usepackage{amssymb}
\usepackage{subcaption}


\title{Projekt z przedmiotu \textit{Podstawowe instrumenty finansowe}}
\author{Aleksy Bałaziński}
\date{\today}

\begin{document}

\maketitle

\section{Wstęp}
Celem projektu było stworzenie modelu, który pozwoli na zrozumienie natury oszczędności emerytalnych oraz zbadanie jego zmienności względem szeregu różnych parametrów.

W projekcie przyjmujemy następujące oznaczenia:
\begin{enumerate}
    \item $\Pi$ -- roczna realna stopa inflacji,
    \item $c$ -- roczna realna stopa zwrotu z obligacji,
    \item $I$ -- roczna realna stopa wzrostu pensji,
    \item $t_0$ -- rok rozpoczęcia odkładnia na emeryturę,
    \item $T$ -- rok przejścia na emeryturę,
    \item $t_d$ -- rok śmierci,
    $p$ -- procent pensji, jaki jest odkładany na emeryturę w każdym miesiącu (funkcja pozostałych parametrów).
\end{enumerate}
Początkowo zakładamy następujące wartości parametrów: $\Pi = 0.02$, $c = 0.03$, $I=0.05$, $t_0 = 0$, $T=40$, $t_d = 60$.
Ze względu na fakt, że wyniki zależą tylko od długości okresu pracy ($T - t_0$) oraz pobierania emerytury ($t_d - T$), dokonano zmiany $t_0 \to t_0 + 25$ (i analogicznie dla pozostałych parametrów czasowych) celem łatwiejszej interpretacji wykresów.

\section{Obliczenie wartości \texorpdfstring{$p$}{p}}
Zakładamy, że jedyna inwestycja jaką mamy do dyspozycji to obligacja ze stopą $c$.
Dodatkowo chcemy, aby pierwsza emerytura była równa ostatniej pensji, a następnie rosła w takim samym tempie jak inflacja.

Niech $r_\Pi$ oznacza miesięczną stopę inflacji, tj. $r_\Pi = (1+\Pi)^{1/12} - 1$.
Analogicznie definiujemy $r_c$ i $r_I$.
Liczba przepracowanych miesięcy to $N = (T-t_0) \cdot 12$, a liczba miesięcy przeżytych na emeryturze to $M=(t_d - T) \cdot 12$.

Jeśli przez $S_0$ oznaczymy początkową pensję, to pensja w miesiącu $k$ wynosić będzie $S_0 (1+r_I)^k$.
Zatem kapitał zgromadzony na potrzeby emerytury przez $N$ miesięcy pracy wyniesie
\begin{equation*}
    W_T = \sum_{k=1}^N pS_0(1+r_I)^{k-1}(1+r_c)^{N-k+1} = pS_0\frac{(1+r_c)^{N+1}}{1+r_I}\sum_{k=1}^N \left( \frac{1+r_I}{1+r_c} \right)^k.
\end{equation*}
Wynika to z faktu, że część pierwszej pensji odłożonej na emeryturę wynosi $p S_0$, co po zakumulowaniu przez stopę $r_c$ przez $N$ miesięcy daje $p S_0 (1+r_c)^N$ itd.

Następnie pobierana jest emerytura.
Pierwsza wypłata $P_0$ wynosi $S_N = S_0 (1+r_I)^N$ i każda kolejna jest powiększona o miesięczną stopę inflacji, a pozostały kapitał w chwili wypłaty $P_0$ jest równy $K_0 = W_T$.
Ogólnie kapitał pozostały tuż przed wypłatą $n$-tej emerytury wynosi $K_n = (K_{n-1} - P_{n-1})(1+r_c)$, a $n$-ta emerytura jest równa $P_n = P_{n-1}(1+r_\Pi) = P_0 (1+r_\Pi)^n$ ($1 \leq  n \leq M$).
Rozpisując wzór na $K_n$ otrzymujemy
\begin{equation}\label{eq:WT}
    K_n = (1+r_c)K_{n-1} - (1+r_c)P_0(1+r_\Pi)^{n-1},
\end{equation}
zatem jest to liniowe równanie różnicowe pierwszego rzędu.
Jego rozwiązanie jest równe
\begin{equation*}
    K_n = (1+r_c)^n K_0 - P_0 \frac{(1+r_c)^{n+1}}{1+r_\Pi} \sum_{k=1}^n \left( \frac{1+r_\Pi}{1+r_c} \right)^k.
\end{equation*}
Optymalne $p$ wyznaczymy z warunku $K_M = 0$ (w chwili śmierci nie zostały nam żadne fundusze).
Otrzymujemy wówczas
\begin{equation*}
    W_T = S_0 (1+r_I)^N\frac{1+r_c}{1+r_\Pi} \sum_{k=1}^M \left( \frac{1+r_\Pi}{1+r_c} \right)^k,
\end{equation*}
a wstawiając $W_T$ z \autoref{eq:WT} i rozwiązując ze względu na $p$, dostajemy
\begin{equation}\label{eq:p-exact-formula}
    p 
    % = \frac{(1+r_I)^N}{1+r_\Pi} \sum_{k=1}^M \left( \frac{1+r_\Pi}{1+r_c} \right)^k 
    % \bigg/
    % \frac{(1+r_c)^N}{1+r_I} \sum_{k=1}^N \left( \frac{1+r_I}{1+r_c} \right)^k
    = \frac{(1+r_I)^{N+1}}{(1+r_\Pi)(1+r_c)^N } \sum_{j=1}^M \left( \frac{1+r_\Pi}{1+r_c} \right)^j 
    \bigg/ \sum_{k=1}^N \left( \frac{1+r_I}{1+r_c} \right)^k.
\end{equation}
Wstawienie wartości parametrów ze wstępu daje wynik $p = 65.1926\%$.
Celem weryfikacji powyższego wyprowadzenia, w programie zaimplementowano również funkcję numerycznego wyznaczenia wartości $p$.
W tym podejściu funkcja \texttt{root\_scalar} z pakietu \texttt{scipy.optimize} jest wykorzystywana do znalezienia $p$, dla którego na koniec życia nie zostają nam żadne fundusze.

Obydwa podejścia dają te same wyniki, co przedstawiono na \autoref{fig:formula-vs-simulation}.
\begin{figure}[htbp]
    \centering
    \includegraphics[width=0.5\linewidth]{figures/formula-vs-simulation.pdf}
    \caption{Kapitał końcowy: porównanie wyników uzyskanych metodą symulacji i ze wzoru}
    \label{fig:formula-vs-simulation}
\end{figure}

\section{Wrażliwość \texorpdfstring{$p$}{p} na parametry modelu}
\subsection{Wpływ pojedynczych parametrów}
W modelu mamy do dyspozycji wiele parametrów.
W tej sekcji zbadamy jak ich zmiany wpływają na wartość $p$, czyli części wypłaty przeznaczanej co miesiąc na emeryturę.

Wyniki analizy przedstawione są na \autoref{fig:main_impact_plots}.

\begin{figure}[htbp]
    \centering
    % --- Row 1 ---
    \begin{subfigure}[b]{0.48\linewidth}
        \centering
        \includegraphics[width=\linewidth]{figures/Impact-of-Salary-Growth-I.pdf}
        \caption{Wpływ wzorstu pensji $I$}
        \label{fig:salary_growth}
    \end{subfigure}
    \hfill
    \begin{subfigure}[b]{0.48\linewidth}
        \centering
        \includegraphics[width=\linewidth]{figures/Impact-of-Bond-Return-c.pdf}
        \caption{Wpływ stopy obligacji $c$}
        \label{fig:bond_return}
    \end{subfigure}
    
    \vspace{0.5cm} % Adds vertical spacing between rows

    % --- Row 2 ---
    \begin{subfigure}[b]{0.48\linewidth}
        \centering
        \includegraphics[width=\linewidth]{figures/Impact-of-Inflation-Pi.pdf}
        \caption{Wpływ inflacji $\Pi$}
        \label{fig:inflation}
    \end{subfigure}
    \hfill
    \begin{subfigure}[b]{0.48\linewidth}
        \centering
        \includegraphics[width=\linewidth]{figures/Impact-of-Retirement-Age-T.pdf}
        \caption{Wpływ wieku $T$ przejścia na emeryturę}
        \label{fig:retirement_age}
    \end{subfigure}

    \vspace{0.5cm}

    % --- Row 3 ---
    \begin{subfigure}[b]{0.48\linewidth}
        \centering
        \includegraphics[width=\linewidth]{figures/Impact-of-Starting-Age-t0.pdf}
        \caption{Wpływ wieku $t_0$ rozpoczęcia oszczędzania na emeryturę}
        \label{fig:starting_age}
    \end{subfigure}
    \hfill
    \begin{subfigure}[b]{0.48\linewidth}
        \centering
        \includegraphics[width=\linewidth]{figures/Impact-of-Age-of-Death-td.pdf}
        \caption{Wpływ wieku śmierci $t_d$}
        \label{fig:age_of_death}
    \end{subfigure}

    % Main Figure Caption
    \caption{Analiza wpływu parametrów modelu na wartość $p$}
    \label{fig:main_impact_plots}
\end{figure}

Z wykresu \autoref{fig:salary_growth} widzimy, że wraz ze wzrostem $I$, procent $p$ jaki musimy odkładać na emeryturę rośnie bardzo szybko.
Może być to w pierwszej chwili dość nieoczekiwane (wraz ze wzrostem pensji, nasze oszczędności emerytalne przecież rosną).
Dzieje się tak z dwu powodów.
Po pierwsze, dla dużego $I$, ostatnia pensja jest wielokrotnie wyższa od początkowej, np. dla $I = 5\%$ po 40 latach $S_N/S_0 \approx 7$ (mamy do czynienia ze wzrostem wykładniczym).
Oznacza to, że pobierana emerytura również będzie bardzo wysoka.
Po drugie, pomimo, że zawsze odkładamy stały procent pensji na emeryturę, nasze zarobki stają się znaczne dopiero po długim okresie czasu (np. dla $I=5\%$, nasza pensja osiągnie dwukrotność początkowej po 14 latach, a czterokrotność po 28 -- pamiętajmy, że $T=40$).

Okazuje się, że zależność $p$ od $r_I$ jest w dobrym przybliżeniu liniowa.
Aby to zobaczyć, spójrzmy na \autoref{eq:p-exact-formula}.
Niech $S_M$ i $S_N$ oznaczają sumy odpowiednio po $j$ (od 1 do $M$) i po $k$ (od 1 do $N$).
Wówczas $S_N = u (u^N - 1)/(u-1) \approx u^{N+1}/(u - 1)$, gdzie $u = (1+r_I) / (1+r_c) > 1$.
Ponadto $u-1 = (r_I - r_c) / (1+r_c)$.
Wstawiając te wyniki do \autoref{eq:p-exact-formula} dostajemy przybliżenie w przypadku, gdy $r_I > r_c$ (dokładne dla dużych $N$)
\begin{equation*}
    p \approx \frac{S_M}{1+r_\Pi}(r_I - r_c).
\end{equation*}

Na \autoref{fig:bond_return} widzimy, że wraz ze wzrostem stopy zwrotu $c$ z obligacji, $p$ maleje.
Jest to oczekiwany rezultat, bowiem większe $c$ pozwala na szybszą akumulację kapitału.
Z \autoref{eq:p-exact-formula} widzimy, że w granicy $c \to \infty$, $p$ dąży do zera.
Faktycznie, dla dużych wartości $r_c$, jedynie pierwsze wyrazy wpływają na sumy $S_N$ i $S_M$.
W takim wypadku mamy $p \propto 1/(1+r_c)^N$.

W efekcie wzrostu stopy inflacji $\Pi$, $p$ rośnie, co widać na \autoref{fig:inflation}.
Jest to powodowane faktem, że wówczas pobierana emerytura również rośnie wykładniczo.


Na \autoref{fig:starting_age} i \autoref{fig:age_of_death} przedstawiono zależność $p$ odpowiednio od wieku rozpoczęcia oszczędzania i od wieku śmierci przy ustalonym wieku $T$ przejścia na emeryturę.
Obrazują one, że zgodnie z oczekiwaniami, odwlekanie w czasie momentu $t_0$ działa na naszą niekorzyść.
Również długi czas pobierania emerytury oznacza konieczność zwiększenia $p$.
Zauważmy jednak, że tempo wzrostu funkcji na wykresie \autoref{fig:age_of_death} wydaje się zmniejszać.
Jest tak ze względu na fakt, że przyjęliśmy $\Pi < c$, a zatem szereg po $j$ w \autoref{eq:p-exact-formula} (dla $M \to \infty$) jest zbieżny.

Na powyższych wykresach dla niektórych wartości parametrów mamy $p > 100\%$.
Oczywiście wynik taki nie ma sensu, jeśli założymy, że pobierana pensja to nasze jedyne źródło dochodu.

\subsection{Interakcje między parametrami}
Celem zbadania interakcji między wybranymi parametrami, narysowano zależność $p$ od tych parametrów, co przedstawiono na \autoref{fig:main_interaction_plots}.
\begin{figure}[htbp]
    \centering
    % --- Row 1 ---
    \begin{subfigure}[b]{0.48\linewidth}
        \centering
        \includegraphics[width=\linewidth]{figures/Interaction-Bond-Yield-c-vs-Inflation-Pi.pdf}
        \caption{Stopa zwrotu z obligacji vs. inflacja}
        \label{fig:interaction_bond_inflation}
    \end{subfigure}
    \hfill
    \begin{subfigure}[b]{0.48\linewidth}
        \centering
        \includegraphics[width=\linewidth]{figures/Interaction-Lifespan-td-vs-Retirement-T.pdf}
        \caption{Czas życia vs. czas przejścia na emeryturę}
        \label{fig:interaction_lifespan_retirement}
    \end{subfigure}
    
    \vspace{0.5cm} % Adds vertical spacing between rows

    % --- Row 2 (Centered) ---
    \begin{subfigure}[b]{0.48\linewidth}
        \centering
        \includegraphics[width=\linewidth]{figures/Interaction-Retirement-T-vs-Starting-Age-t0.pdf}
        \caption{Czas przejścia na emeryturę vs. czas rozpoczęcia odkładnia}
        \label{fig:interaction_retirement_starting}
    \end{subfigure}

    % Main Figure Caption
    \caption{Zależność $p$ od dwu parametrów}
    \label{fig:main_interaction_plots}
\end{figure}
Wykres na \autoref{fig:interaction_bond_inflation} ilustruje, że celem zachowania stałego $p$, inflacja i stopa zwrotu z obligacji muszą zmieniać się wspólnie (razem rosnąć lub maleć).
Dla przykładu przy stałym $p=55\%$, gdy $r_c$ zmienia się z 3\% na 5\%, $r_\Pi$ zmienia się z 0 na 6\%.

Na \autoref{fig:interaction_lifespan_retirement} widzimy, że przy stałym $p$, wiek przejścia na emeryturę i czas śmierci $t_d$ są dodatnio skorelowane.
Jeśli $p=60\%$ i $t_d$ zmienia się z 70 na 90, to $T$ zmienia się z 55 na 70.

Z \autoref{fig:interaction_retirement_starting} wynika, że $p$ jest utrzymane na stałym poziomie, gdy zmianie $T$ towarzyszy zmiana $t_0$ w tym samym kierunku.
Na przykład, przy poziomie 60\%, gdy $T$ zmienia się z 65 na 67.5, $t_0$ zmienia się z 20 na 32.
Uwidacznia to, że niewielka zmiana $T$ wymaga znacznej zmiany $t_0$.

\section{Zgodność modelu z rzeczywistością}
Niektóre z założeń modelu są nierealistyczne i podważają jego stosowalność do planowania emerytalnego:
\begin{itemize}
    \item Zakładamy, że stopa inflacji $\Pi$ i stopa zwrotu z obligacji $c$ pozostają stałe przez 60 lat.
    W praktyce te parametry są zmienne.
    Również założenie, że wzrost pensji $I$ jest stały przez wszystkie 40 lat pracy jest ryzykowne.
    Aby uczynić model bliższym rzeczywistości, należałoby wprowadzić losowość przy ustalaniu wartości tych parametrów na dany rok.
    
    Dodatkowo model zakłada, że będziemy pracować przez cały okres od $t_0$ do $T$.
    Nie bierze on w ogóle pod uwagę możliwości przerw w zatrudnieniu.
    

    Bardziej realistyczny model powinien również uwzględnić możliwość niewypłacalności obligacji, na których opieramy nasze oszczędności.

    Te dwie kwestie również powinny być modelowane w sposób stochastyczny.
    \item W naszym modelu $t_d$ (czas śmierci) jest ustalony z góry.
    O ile wcześniejsza śmierć nie stanowi problemu (przynajmniej z rozważanej perspektywy pobierania emerytury), tak w przypadku dłuższego niż planowany czasu życia, ryzykujemy bankructwem na emeryturze.
    Tutaj ponownie dodanie czynnika losowego uczyniłoby model lepszym.
    \item Zakładamy, że nasza emerytura musi być równa ostatniej pensji (i zwiększa się następnie o stopę inflacji).
    Można uznać to za zdecydowanie zbyt wygórowane oczekiwania i ustalić mniejszy poziom pobieranej emerytury.
    \item Jedną z możliwości dalszego rozwinięcia modelu jest również uwzględnienie opłat i prowizji związanych z prowadzeniem rachunku maklerskiego, na którym trzymamy nasze obligacje.
\end{itemize}
Model powinien również uwzględniać możliwość inwestowania w instrumenty inne niż tylko obligacje.
Tą kwestią zajmiemy się w kolejnym paragrafie.

\section{Oszczędzanie z wykorzystaniem indeksu giełdowego}
W tej części zakładamy, że mamy do dyspozycji drugie aktywo w postaci indeksu giełdowego.
Jest on scharakteryzowany przez trzy parametry $q_u$ (roczna realna stopa zwrotu w przypadku, gdy indeks daje dodatnie zwroty z inwestycji), $q_d$ (roczna realna stopa zwrotu w przypadku, gdy indeks daje ujemne zwroty z inwestycji) oraz $p_u$ (prawdopodobieństwo uzyskania dodatnich zwrotów z funduszu).
Oprócz tego definiujemy $p_d = 1 - p_u$.

\subsection{Estymacja parametrów}
Celem określenia realistycznych wartości wymienionych parametrów przeanalizowano dane historyczne dla indeksu S\&P 500.
Wykorzystano do tego zbiór danych opublikowany przez New York University dostępny pod adresem \url{https://pages.stern.nyu.edu/~adamodar/New_Home_Page/datafile/histretSP.html}.
Podane tam wartości nie uwzględniają inflacji (figurują tam jako \textit{Annual Returns on Investment}), zatem wykorzystano wzór 
\begin{equation*}
    r = \frac{1 + i}{1 + \Pi} - 1
\end{equation*}
do zamiany stóp nominalnych $i$ na rzeczywiste $r$.
Wartości stopy inflacji zostały zaczerpnięte ze strony Banku Światowego (\url{https://data.worldbank.org/indicator/FP.CPI.TOTL.ZG?locations=US}).
Następnie, celem wyznaczenia $p_u$ wzięto liczbę lat z dodatnią realną stopą zwrotu przez całkowitą rozważaną liczbę lat (za punkt startowy wzięto rok 1980).
Wartości $q_u$ i $q_d$ zostały oszacowane jako średnia realna stopa zwrotu liczona dla odpowiednio lat z dodatnimi zwrotami i ujemnymi.
Tą metodą uzyskano wartości $p_u=0.7609$, $q_u = 0.166$ i $q_d = -0.1429$.
Daje to średnią realną stopę równą $9.2\%$ i odchylenie standardowe równe $13.2\%$.

\subsection{Strategia \textit{glide path}}
W tej części projektu przyjęto model opierający się na tzw. strategii \textit{glide path}, tj. strategii której celem jest stopniowa zmiana struktury portfela z bardziej agresywnej (akcje) na bardziej konserwatywną (obligacje) w miarę zbliżania się do określonej daty docelowej.
Dodatkowo przyjęto, że udział akcji w portfelu będzie zmniejszał się liniowo aż do czasu $T$.
Od tego momentu skład portfela nie zmienia się.
Zatem w modelu mamy dwa dodatkowe parametry: $w_s$ i $w_e$, obydwa między 0 i 1, które określają względny udział akcji w portfelu odpowiednio na początku i na końcu okresu zarobkowania (w czasie od $t_0$ do $T$), a nasza \textit{strategia inwestycyjna} określona jest w sumie przez trzy liczby: $p$, $w_s$ i $w_e$.

Dodanie do modelu akcji sprawia, że wprowadzony zostaje czynnik losowy.
W zaimplementowanej symulacji na początku każdego roku z prawdopodobieństwem $p_u$ podejmowana jest decyzja, czy dany rok przyniesie zyski na akcjach (stopa $q_u > 0$), czy straty (stopa $q_d < 0$).
Strategię uznajemy za udaną jeśli w 95\% przeprowadzonych prób w momencie śmierci nasze fundusze były nieujemne.

Celem znalezienia optymalnego $p$ przetestowano wartości 
\begin{equation}\label{eq:ws-we-values}
    w_s \in \{0, .7, .8, .9, 1\}, \quad w_e \in \{0, .2, .3, .4, .5\}.
\end{equation}
Dla każdej pary $(w_s, w_e)$ wyznaczono optymalną wartość $p$ (tj. najmniejszą, która prowadzi do sukcesu w 95\% przypadków).

Dla rozważanych wartości parametrów, optymalną okazała się strategia z $w_s = 1$, $w_e = 0.5$ i $p=44.64\%$.

Rozkład wartości pozostałych po śmierci funduszy dla wyznaczonych parametrów przedstawia \autoref{fig:dist-final-wealth}.
\begin{figure}[htbp]
    \centering
    \includegraphics[width=0.75\linewidth]{figures/dist-final-wealth.pdf}
    \caption{Rozkład wartości pozostałych po śmierci funduszy dla strategii z $w_s = 1$, $w_e = 0.5$ i $p=44.64\%$ (pensja początkowa równa \$1000).}
    \label{fig:dist-final-wealth}
\end{figure}
Do jego wygenerowania przejęto, że pierwsza pensja (w momencie $t_0$) była równa \$1000.
Uzyskany histogram ma cechy charakterystyczne dla rozkładu log-nromalnego, który pojawia się w kontekście analizy cen akcji w modelu dwumianowym.
W szczególności zauważyć można długi prawy ogon, który odpowiada rzadkim sytuacjom, w którym akcje przynoszą ogromne zyski.
Jednocześnie w większości przypadków otrzymujemy ,,umiarkowane zyski'' z gry na giełdzie.

Celem sprawdzenia jak dobór różnych wartości $p$ wpływa na szansę uniknięcia bankructwa na emeryturze, sprawdzono w jakim procencie przypadków znaleziona strategia z $w_s = 1$ i $w_e = 0.5$ daje nieujemny kapitał w chwili śmierci.
Oczywiście spodziewamy się, że wraz ze wzrostem $p$ nasze szanse również rosną.
Wynik eksperymentu jest przedstawiony na \autoref{fig:p-vs-succ-rate}.
\begin{figure}[htbp]
    \centering
    \includegraphics[width=0.5\linewidth]{figures/p-vs-succsses-rate.pdf}
    \caption{Wpływ $p$ na szanse uniknięcia bankructwa na emeryturze.}
    \label{fig:p-vs-succ-rate}
\end{figure}
Dodatkowo na rysunku narysowano analogiczny wykres dla strategii rozważanej w poprzednich paragrafach (portfolio składające się wyłącznie z obligacji).
Widzimy, że portfel z akcjami i obligacjami daje nam większą ,,elastyczność'' w porównaniu do portfela zawierającego tylko obligacje.
W tym drugim przypadku funkcja sukcesu jest skokowa (przeskok z wartości 0 na 1 występuje dla obliczonego wcześniej $p\approx 65\%$).
Na podstawie tego wykresu możemy zadecydować jakie $p$ jest dla nas najlepsze w zależności od naszej skłonności do poodejmowania ryzyka (oczywiście w kontekście emerytalnym nie powinna ona być zbyt duża).

\subsection{Analiza wrażliwości na dwa parametry}
Dla zaproponowanej w poprzednim paragrafie strategii \textit{glide path}, zbadano jaki wpływ mają parametry modelu na optymalne wartości $p$, $w_s$ i $w_e$.
Podobnie jak poprzednio, ograniczono się do testowania kombinacji wartości parametrów $w_e$ i $w_s$ określonych w (\ref{eq:ws-we-values}).
Wyniki eksperymentu przedstawia \autoref{fig:main_optimal_savings_plots}.
\begin{figure}[htbp]
    \centering
    % --- Row 1 ---
    \begin{subfigure}[b]{0.49\linewidth}
        \centering
        \includegraphics[width=\linewidth]{figures/Optimal-Savings-Rate-p-and-Glide-Path-vs-c-and-Pi.pdf}
        \caption{Optymalna stopa oszczędności vs. c i inflacja $\Pi$}
        \label{fig:optimal_savings_c_pi}
    \end{subfigure}
    \hfill
    \begin{subfigure}[b]{0.49\linewidth}
        \centering
        \includegraphics[width=\linewidth]{figures/Optimal-Savings-Rate-p-and-Glide-Path-vs-td-and-T.pdf}
        \caption{Optymalna stopa oszczędności vs. $t_d$ i $T$}
        \label{fig:optimal_savings_td_T}
    \end{subfigure}
    
    \vspace{0.5cm} % Adds vertical spacing between rows

    % --- Row 2 (Centered) ---
    \begin{subfigure}[b]{0.49\linewidth}
        \centering
        \includegraphics[width=\linewidth]{figures/Optimal-Savings-Rate-p-and-Glide-Path-vs-T-and-t0.pdf}
        \caption{Optymalna stopa oszczędności vs. $T$ i $t_0$}
        \label{fig:optimal_savings_T_t0}
    \end{subfigure}

    % Main Figure Caption
    \caption{Optymalna stopa oszczędności $p$ oraz ścieżka alokacji (Glide Path) w zależności od wybranych parametrów}
    \label{fig:main_optimal_savings_plots}
\end{figure}
Przedstawione na nim mapy ciepła wyglądają analogicznie do tych przedstawionych na \autoref{fig:main_interaction_plots} z tą różnicą, że w tym przypadku ograniczono się do obliczenia wartości jedynie w dziewięciu punktach z powodów długiego czasu trwania symulacji.

Z \autoref{fig:optimal_savings_c_pi} możemy wywnioskować, że gdy stopa zwrotu z obligacji jest duża a inflacja jest na niskim poziomie, możemy zastosować mniej agresywną strategię.
Na przykład dla $c=5\%$ i $\Pi=2.5\%$, optymalną strategią okazuje się rozpoczęcie odkładania na emeryturę z portfelem zawierającym 70\% akcji, a na emeryturze korzystanie z portfela zawierającego 30\% akcji.
W takiej sytuacji $p$ jest również zdecydowanie mniejsze niż 65\% uzyskane dla parametrów bazowych i wynosi 37.4\%.
Z drugiej strony, jeśli inflacja jest bardzo duża względem $c$, zarówno $w_s$ jak i $w_e$ powinny być znaczne, np. dla $c=1\%$ i $\Pi = 4\%$, $w_s = 1$ i $w_e=0.5$ (są to maksymalne rozważane wartości tych parametrów).

Sytuacja na \autoref{fig:optimal_savings_td_T} jest zupełnie analogiczna do tej przedstawionej na \autoref{fig:interaction_lifespan_retirement}, z tą różnicą, że nawet w przypadku, gdy $T = t_d = 70$ (tzn., że nie zdążymy nacieszyć się emeryturą), $p > 0$.
Jest to związane z faktem, że w momencie zakończenia pracy, na skutek niekorzystnych ruchów na giełdzie, nasz odłożony kapitał może być ujemny.
Dodatkowo, mimo iż metoda glide path pozwala na zmniejszenie $p$, dla zbyt dużych okresów emerytury (np. $T$ = 50, $t_d = 100$) nie pozwala ona na znalezienie satysfakcjonujących rozwiązań.

Porównanie \autoref{fig:optimal_savings_T_t0} z \autoref{fig:interaction_retirement_starting} pokazuje, że zastosowanie agresywnej strategii ($w_s = 0.5$, $w_e = 1$) pozwala na uzyskanie sukcesu w sytuacjach, gdzie portfel oparty o same obligacje był niewystarczający (np. pary $(T, t_0) = (50, 20), (60, 35)$).

\subsection{Analiza wrażliwości na jeden parametr}
Celem dokładniejszego poznania jak różne parametry wpływają na optymalne wartości $p$, $w_s$ i $s_e$, narysowano wykresy analogiczne do tych przedstawionych na \autoref{fig:main_impact_plots}.
Uzyskane wykresy przedstawia \autoref{fig:main_glide_path_impact_plots}.
\begin{figure}[htbp]
    \centering
    % --- Row 1 ---
    \begin{subfigure}[b]{0.48\linewidth}
        \centering
        \includegraphics[width=\linewidth]{figures/Impact-of-Salary-Growth-I-Glide-Path.pdf}
        \caption{Wpływ wzrostu pensji $I$ na ścieżkę alokacji}
        \label{fig:glide_path_salary_growth}
    \end{subfigure}
    \hfill
    \begin{subfigure}[b]{0.48\linewidth}
        \centering
        \includegraphics[width=\linewidth]{figures/Impact-of-Bond-Return-c-Glide-Path.pdf}
        \caption{Wpływ stopy obligacji $c$ na ścieżkę alokacji}
        \label{fig:glide_path_bond_return}
    \end{subfigure}
    
    \vspace{0.5cm} % Adds vertical spacing between rows

    % --- Row 2 ---
    \begin{subfigure}[b]{0.48\linewidth}
        \centering
        \includegraphics[width=\linewidth]{figures/Impact-of-Inflation-Pi-Glide-Path.pdf}
        \caption{Wpływ inflacji $\Pi$ na ścieżkę alokacji}
        \label{fig:glide_path_inflation}
    \end{subfigure}
    \hfill
    \begin{subfigure}[b]{0.48\linewidth}
        \centering
        \includegraphics[width=\linewidth]{figures/Impact-of-Retirement-Age-T-Glide-Path.pdf}
        \caption{Wpływ wieku emerytalnego $T$ na ścieżkę alokacji}
        \label{fig:glide_path_retirement_age}
    \end{subfigure}

    \vspace{0.5cm}

    % --- Row 3 ---
    \begin{subfigure}[b]{0.48\linewidth}
        \centering
        \includegraphics[width=\linewidth]{figures/Impact-of-Starting-Age-t0-Glide-Path.pdf}
        \caption{Wpływ wieku $t_0$ rozpoczęcia na ścieżkę alokacji}
        \label{fig:glide_path_starting_age}
    \end{subfigure}
    \hfill
    \begin{subfigure}[b]{0.48\linewidth}
        \centering
        \includegraphics[width=\linewidth]{figures/Optimized-Impact-of-Age-of-Death-td-Glide-Path.pdf}
        \caption{Wpływ wieku śmierci $t_d$ na ścieżkę alokacji}
        \label{fig:glide_path_age_of_death}
    \end{subfigure}

    % Main Figure Caption
    \caption{Analiza wpływu parametrów modelu na ścieżkę alokacji (Glide Path)}
    \label{fig:main_glide_path_impact_plots}
\end{figure}
Wnioski, które można z nich wyciągnąć są podobne do tych opisanych przy okazji omawiania \autoref{fig:main_impact_plots}, zatem nie będą tutaj powtarzane.

Na uwagę zasługuje fakt, że w przypadku strategii glide path uzyskane wartości $p$ są zawsze mniejsze od odpowiadających im wartości w przypadku portfela składającego się wyłącznie z obligacji.
Również możemy stwierdzić, że dla wykorzystanej w projekcie kombinacji parametrów bazowych, jedynie zwiększenie stopy zwrotu z obligacji $c$ miało wpływ na optymalne wartości $w_s$ i $w_e$, co widać na \autoref{fig:glide_path_bond_return} (dla $c > 4\%$, ich optymalne wartości zmieniły się z $(1, 0.5)$ na $(0.7, 0.2)$).

\section{Podsumowanie i wnioski}
Pomimo iż modele, które zostały rozważone w raporcie były znacznie uproszczone, pozwoliły one na wyciągnięcie kilku ważnych wniosków dotyczących planowania oszczędności emerytalnych.
Do najważniejszych konkluzji należą:
\begin{itemize}
    \item \textbf{Kluczowa rola czasu (efekt procentu składanego):} Wczesne rozpoczęcie odkładania na emeryturę ($t_0$) radykalnie zmniejsza wymaganą miesięczną stopę oszczędności ($p$). Każdy rok zwłoki skutkuje drastycznym, nieliniowym wzrostem obciążenia bieżącej pensji w późniejszych latach.
    \item \textbf{Paradoks szybkiego wzrostu wynagrodzenia:} Wysoka stopa wzrostu pensji ($I$) -- wbrew intuicji -- zmusza do odkładania większego odsetka dochodów ($p$).
    Wynika to z przyjętego założenia, że wypłacana emerytura ma być równa co najmniej ostatniej pensji.
    Przy wykładniczym wzroście płac pod koniec kariery, kapitał zgromadzony z początkowych, niskich składek staje się relatywnie marginalny w stosunku do oczekiwanej, wysokiej emerytury.
    \item \textbf{Wpływ czynników makroekonomicznych:} Model wykazuje ogromną wrażliwość na realną stopę zwrotu z obligacji ($c$) oraz inflację ($\Pi$).
    Aby utrzymać stałą stopę oszczędności, parametry te muszą rosnąć lub maleć korelując ze sobą –- w przeciwnym wypadku wysoka inflacja bez adekwatnego wzrostu rentowności bezpiecznych aktywów prowadzi do matematycznej niemożliwości sfinansowania emerytury ($p > 100\%$).
    \item \textbf{Przewaga strategii (\textit{glide path}):} Wprowadzenie do portfela komponentu akcyjnego (indeksu giełdowego) znacząco poprawia efektywność modelu.
    Optymalna strategia oparta na liniowej redukcji ryzyka ($w_s = 1 \to w_e = 0.5$) pozwoliła na obniżenie wymaganej stopy oszczędności z bazowych $65.19\%$ (dla samych obligacji) do znacznie bardziej realistycznego poziomu $44.64\%$, przy zachowaniu $95\%$ prawdopodobieństwa sukcesu.
    \item \textbf{Cena bezpieczeństwa a elastyczność struktury:} Portfel oparty wyłącznie o obligacje cechuje się zero-jedynkową funkcją sukcesu (sztywny próg rentowności).
    Włączenie instrumentów o charakterze stochastycznym (akcje) wygładza tę funkcję, dając inwestorowi elastyczność i możliwość świadomego dopasowania stopy oszczędności do własnej awersji do ryzyka, co odzwierciedla rozkład lognormalny końcowych funduszy z długim, korzystnym prawym ogonem.
\end{itemize}

W ramach dalszego rozwoju projektu, naturalnym krokiem byłoby odejście od modeli deterministycznych w sekcji obligacji na rzecz symulacji stochastycznych (np. procesów CIR lub Vasiceka dla stóp procentowych) oraz modelowanie czasu przeżycia ($t_d$) za pomocą tablic trwania życia.
Pozwoliłoby to na jeszcze lepsze odzwierciedlenie realiów rynkowych i demograficznych.

\end{document}

